Since @@M0@@ is @@M1@@-regular, @@M2@@ by 1.4; hence @@M3@@. Now @@M4@@ by 1.8 (i), since @@M5@@ is a @@M6@@-sheaf by 2.5. Consequently,
\sgaStableTarget{sga6.e13.eq.2-6-2}@@M7@@
Finally, 2.7.2 gives (ii), and 2.7.1 and 2.7.2 give (i).

\sgaStableTarget{sga6.e13.lem.2-8}\textbf{Lemma 2.8.} \emph{We suppose that @@M8@@ and @@M9@@. Let @@M10@@ and @@M11@@. Then}:
@@M12@@

Indeed, let @@M13@@; then, by comparing polynomials, it follows that @@M14@@, that @@M15@@ and that @@M16@@. Let @@M17@@; then @@M18@@ and, by 2.3, @@M19@@, which equals @@M20@@. Thus @@M21@@, and since @@M22@@, we have @@M23@@ by 2.7 (ii); this gives the assertion.

\sgaStableTarget{sga6.e13.lem.2-9}\textbf{Lemma 2.9.} \emph{We suppose that @@M24@@, that @@M25@@ and that @@M26@@ is normal. Then}:
\begin{enumerate}[label=(\roman*)]
\item \emph{@@M27@@ for @@M28@@,}
\item \emph{@@M29@@ for every @@M30@@,}
\item \emph{@@M31@@ for @@M32@@,}
\end{enumerate}
\emph{where @@M33@@ and @@M34@@ are universal polynomials with coefficients in @@M35@@, which depend on neither @@M36@@ nor @@M37@@.}

